\documentclass[conf]{new-aiaa} % Use [journal] for AIAA journal style

%\usepackage[utf8]{inputenc}

%graphics packages
\usepackage[]{graphicx}
\graphicspath{{./figures/}}

% ------------------ tikz ------------
%%tikz
\RequirePackage{pgf,tikz}
\usepackage{pgfplots}
\pgfplotsset{compat=newest}
\usetikzlibrary{arrows.meta}
\usetikzlibrary{arrows}
\usetikzlibrary{shapes.geometric}
\usetikzlibrary{shapes.misc}
\usetikzlibrary{backgrounds}
\usetikzlibrary{bending}
\usetikzlibrary{shadows}
\usetikzlibrary{patterns}
\usetikzlibrary{patterns.meta}
\usetikzlibrary{intersections}
\usepgfplotslibrary{patchplots}
\usepgfplotslibrary{fillbetween}
\usepgfplotslibrary{groupplots}
\usepgfplotslibrary{polar}
\usepgfplotslibrary{smithchart}
\usepgfplotslibrary{statistics}
\usepgfplotslibrary{dateplot}
\usepgfplotslibrary{ternary}
\pgfplotsset{%
	layers/standard/.define layer set={%
		background,axis background,axis grid,axis ticks,axis lines,axis tick labels,pre main,main,axis descriptions,axis foreground%
	}{grid style= {/pgfplots/on layer=axis grid},%
		tick style= {/pgfplots/on layer=axis ticks},%
		axis line style= {/pgfplots/on layer=axis lines},%
		label style= {/pgfplots/on layer=axis descriptions},%
		legend style= {/pgfplots/on layer=axis descriptions},%
		title style= {/pgfplots/on layer=axis descriptions},%
		colorbar style= {/pgfplots/on layer=axis descriptions},%
		ticklabel style= {/pgfplots/on layer=axis tick labels},%
		axis background@ style={/pgfplots/on layer=axis background},%
		3d box foreground style={/pgfplots/on layer=axis foreground},%
	},
}
% new style at automates partial ellipse
\tikzset{
    partial ellipse/.style args={#1:#2:#3}{
        insert path={+ (#1:#3) arc (#1:#2:#3)}
    }
}

\usepackage{caption}
\usepackage{subcaption}
\usepackage{placeins}
\usepackage{wrapfig}

%tables
\usepackage{array}
\newcommand{\PreserveBackslash}[1]{\let\temp=\\#1\let\\=\temp}
\newcolumntype{C}[1]{>{\PreserveBackslash\centering}m{#1}}
\newcolumntype{R}[1]{>{\PreserveBackslash\raggedleft}m{#1}}
\newcolumntype{L}[1]{>{\PreserveBackslash\raggedright}m{#1}}

%mathpackages
\usepackage{amsmath}
%\usepackage{siunitx}
%\usepackage{mathrsfs}
%\newcommand{\vect}{\overset{\rightharpoonup}}
%decrease overset height
% \makeatletter
% \newcommand{\oset}[3][0.25ex]{%
% 	\mathrel{\mathop{#3}\limits^{
% 			\vbox to#1{\kern-0.5\ex@
% 				\hbox{$\scriptstyle#2$}\vss}}}}
% \makeatother

%%Vector arrow over variable
%\newcommand{\vect}[1]{%
%	\oset{\rightharpoonup}{#1}}

\RequirePackage{bm}
\newcommand{\vect}[1]{\bm{#1}}

\newcommand{\mat}[1]{%
	\mathbf{#1}}
\usepackage{nicefrac}
\usepackage{cancel}

%referencing packages
\usepackage{hyperref}
\usepackage[]{cleveref}

% ---------- colors -------------
\RequirePackage{xcolor}
\definecolor{navy}{HTML}{002E5D}
\definecolor{royal}{HTML}{005CAB}	% royal that matches BYU Engineering logo
\definecolor{darkgray}{HTML}{141414}
\definecolor{mediumgray}{HTML}{666666}	% medium gray definition from A. Ning
\definecolor{black}{HTML}{111111}
\definecolor{primary}{HTML}{005CAB}
\definecolor{secondary}{HTML}{c05367}
\definecolor{tertiary}{HTML}{8fa651}
\definecolor{plotsgray}{HTML}{808080}


% -------------- easy coloring of things ----------------
\newcommand{\navy}[1]{{\color{navy}#1}}
\newcommand{\primary}[1]{{\color{primary}#1}}
\newcommand{\secondary}[1]{{\color{secondary}#1}}
\newcommand{\tertiary}[1]{{\color{tertiary}#1}}
\newcommand{\gray}[1]{{\color{gray}#1}}

% Allow same footnotes
\newcommand*\samethanks[1][\value{footnote}]{\footnotemark[#1]}

% make \noindent where easier
\newcommand{\where}{\noindent where }
 % your preamble.tex file with packages

\begin{document}
\sloppy % helps reduce overfull hboxes

\title{Notes on the Fast Multipole Method, Vortex Particle Method, and Neural Approximations}
\author{Nathan Lehnhof}
\date{\today}
\maketitle

\tableofcontents

% ======================================================
\section{Mathematical and Physical Foundations}

\subsection{Vector Calculus Review}
Divergence: $\nabla \cdot \mathbf{F}$ measures the net flow out of a point.  
Gradient: $\nabla \phi$ gives the direction of steepest increase.  
Curl: $\nabla \times \mathbf{F}$ represents local rotation.

\subsection{Integral Theorems (Gauss, Stokes, Green)}
Gauss's theorem links divergence to flux through a surface.  
Stokes's theorem relates curl to circulation along a boundary.  
Green's identities provide integral relations for harmonic functions.

\subsection{Inner and Outer Products}
The dot product measures alignment; the cross product encodes area and orientation.

\subsection{Potential Theory}
Laplace’s equation $\nabla^2 \phi = 0$ governs harmonic functions.  
Fundamental solutions and Green’s functions are tools to construct solutions.

\subsection{Numerical Analysis Essentials}
Error, stability, and convergence define the quality of algorithms.  
Interpolation and quadrature enable discrete approximations of continuous functions.  
Computational complexity determines feasibility of large‐scale simulations.

% ======================================================
\section{Fluid Dynamics Foundations}

\subsection{Vorticity and Velocity Formulations}
The vorticity field $\boldsymbol{\omega} = \nabla \times \mathbf{u}$ captures local rotation.  

\subsection{Biot--Savart Law}
Velocity induced by a vortex filament can be written as
\[
\mathbf{u}(\mathbf{x}) = \frac{\Gamma}{4 \pi} \int \frac{\mathrm{d}\boldsymbol{\ell} \times (\mathbf{x}-\mathbf{x}')}{\|\mathbf{x}-\mathbf{x}'\|^3}.
\]

\subsection{Vortex Particle Method}
Vorticity is discretized into particles with circulation strength.  
Particles advect with the flow and interact via smoothing kernels.

% ======================================================
\section{Fast Multipole Method (FMM)}

\subsection{Motivation: Reducing Complexity}
Naive $O(N^2)$ interactions become $O(N)$ or $O(N \log N)$ with FMM.

\subsection{Multipole Expansions}
Laplace and Helmholtz kernels admit far‐field expansions, reducing interaction cost.

\subsection{Hierarchical Tree Structures}
Quadtrees (2D) and octrees (3D) partition space into boxes, grouping particles adaptively.

\subsection{Translation Operators}
\begin{itemize}
  \item M2M: Multipole to Multipole
  \item M2L: Multipole to Local
  \item L2L: Local to Local
\end{itemize}

\subsection{Complexity Analysis and Error Control}
FMM balances speed with controllable approximation error.

% ======================================================
\section{FMM for the Vortex Particle Method}

\subsection{Coupling Biot--Savart with FMM}
FMM accelerates evaluation of induced velocities from vortex particles.

\subsection{Near‐field vs. Far‐field}
Direct summation is used for nearby particles; multipole expansions handle distant ones.

\subsection{Adaptive Trees}
Nonuniform particle distributions require adaptive octrees for efficiency.

\subsection{Accuracy and Kernel Regularization}
Regularized kernels prevent singular behavior in particle interactions.

% ======================================================
\section{Neural Network Surrogates}

\subsection{Problem Framing}
Inputs: relative position and circulation.  
Outputs: induced velocity vectors.

\subsection{Network Architectures}
Explored options include multilayer perceptrons (MLPs), graph neural networks (GNNs), and equivariant neural networks that respect rotational symmetry.

\subsection{Training Data}
Ground truth is generated from Biot--Savart computations.  
Supervised learning minimizes prediction error relative to true velocities.

\subsection{Deployment Strategies}
Neural surrogates replace kernel evaluations in FMM.  
Hybrid methods combine classical and neural approaches.  
Physical constraints are enforced to preserve conservation laws.

% ======================================================
\section{Advanced Topics and Implementation}

\subsection{Error Control and Stability}
Numerical stability is essential in long‐time vortex simulations.

\subsection{Parallelization and GPU Acceleration}
Tree traversal and kernel evaluation map well to GPU architectures.

\subsection{Real‐Time Solver Design}
Combining FMM and neural surrogates enables real‐time vortex simulations.

\subsection{Case Studies}
Applications include vortex rings, wakes, and jet flows.

% ======================================================
\section{References and Resources}

% \nocite{*}
% \bibliographystyle{new-aiaa}
% \bibliography{references}

\end{document}
